\scanpage{67}\mspage{12}
\deleted{plus celle d'une} de ce torseur, \deleted{qui} \inserted{Ce torseur} peut aussi s'interpréter \deleted{comme le torseur des sections} \inserted{ainsi : la structure d'extension (39)} définit la structure d'extension
\diagram{\begin{tikzcd}[ampersand replacement=\&,column sep=small,row sep=0.3em]
\& \& \underline{\mathrm{End}}(M_0) \& \& \underline{\mathrm{End}}(M_0) \\
\& \& \rotatebox{90}{$\simeq$} \& \& \rotatebox{90}{$\simeq$} \\
\text{(43)} \& \unsure{0'|0} \arrow[r] \& V_{\underline{J}}(M_0) \arrow[r] \& V_{\underline J}(M) \arrow[r] \& V_J(M_0) \arrow[r] \& 1
\end{tikzcd}}
où pour \unsure{tt|le|$A$} schéma \inserted{(ou faisceau)} en groupes comm. $A$ sur $S$, on pose
\[
V_{\underline J}(A)\simeq \mathrm{Ker}\big(A^{\underline J}\xrightarrow{\mathrm{somme}} A\big) \tag*{(44)}
\]
(faisceau loc isom. à $A^2$), et $V_{\underline J}$ est un foncteur exact en $A$, commutant bien sûr au chgt. de bas. Pour $A$ annulé par 2, on a un isom
\[
V_{\underline J}(A)\simeq\mathrm{Hom}(M_0,A)\qquad\text{si }2\,\mathrm{id}_A=0 \tag*{(45)}
\]
dû au fait que l'on a
\[
M_0\simeq\mathbb F_2^{\underline J}/\underbrace{\mathbb F_2}_{\mathrm{diag}} \tag*{(46)}
\]

Donc à une section \deleted{\unsure{de|d}} canonique partout non nulle de $V_J(M_0)$, correspondant à l'application d'inclusion $\underline J\simeq M_0^\ast\hookrightarrow M_0$, ou encore à la section \deleted{\unsure{``iden|``idem|évidente}} $\mathrm{id}_{M_0}$ de $\underline{\mathrm{End}}(M_0)$, dont l'image inverse dans (43) est le torseur cherché.
% archive sheet number: 66
